\bookchapter{الفصل 44: تطبيق عملي من 0 إلى 1 — الدورة الكاملة للمشروع من البداية إلى النهاية}

\chapterintro{انتهى وقت النظرية، وحان وقت العمل. هذا الفصل ليس درسا بسيطا من نوع “انسخ واكتب”، بل منهجية عملية لـ“كيف تفكر، كيف تسأل، كيف تتحقق”. سنعرض من خلال ثلاثة مشاريع بدرجات صعوبة مختلفة عملية التفكير الكاملة من المتطلبات إلى المنتج النهائي. وأي أداة تستخدمها تنطبق عليها هذه المنهجية.}

\booksection{أولاً: المشروع الأول — موقع أدوات شخصي فائق البساطة (للمبتدئين)}

\booksubsection{التفكير في المتطلبات}

الهدف: إنشاء صفحة ويب لموقع أدوات شخصية تعرض مداخل 3-4 أدوات صغيرة عملية. هذا مشروع واجهة أمامية خالص لا يحتاج خادما خلفيا ولا قاعدة بيانات. المهم ليس مدى تعقيد الأدوات نفسها، بل تجربة الدورة الكاملة “من المتطلبات إلى المنتج”.

وفق عناصر المتطلبات الخمسة: مَن (Who) — استخدام شخصي. ماذا (What) — صفحة ويب تعرض مداخل الأدوات. لماذا (Why) — جمع الأدوات كثيرة الاستخدام في صفحة واحدة لتسهيل الوصول. كيف (How) — ملف واحد HTML+CSS+JS دون أطر عمل. القيد (Constraint) — الإنجاز خلال 30 دقيقة دون حاجة إلى النشر.

\booksubsection{إصدار التعليمات}

استنادا إلى مبدأ Think Before Coding (فكر قبل أن تبرمج)، ينبغي أن تتضمن تعليماتك معلومات كافية ليفهم الذكاء الاصطناعي الهدف، دون تحديد كل تفاصيل التنفيذ:

\begin{lstlisting}
أنشئ صفحة ويب شخصية للأدوات وفق هذه المتطلبات:
1. ملف HTML واحد (CSS و JS مضمّنان)، دون أُطر عمل
2. سمة داكنة بأسلوب عصري بسيط
3. اعرض 4 بطاقات للدخول إلى الأدوات:
   - أداة تلخيص النصوص
   - عدّاد الكلمات
   - منتقي الألوان
   - مولّد رموز QR
4. لكل بطاقة أيقونة وعنوان ووصف قصير
5. للبطاقات تأثير حركي عند التحويم
6. تخطيط متجاوب يعمل على الهاتف
7. شريط عنوان في الأعلى
أعطني مباشرة ملف HTML الكامل.
\end{lstlisting}

\booksubsection{القبول والتكرار}

بعد أن يعطيك الذكاء الاصطناعي الكود، اقبل العمل وفق فحوص التحقق الثلاثية. الطبقة الأولى، التحقق الوظيفي: هل يفتح الملف في المتصفح؟ هل ظهرت البطاقات الأربع كلها؟ هل تعمل حركات التحويم (Hover) بشكل سليم؟ الطبقة الثانية، التحقق الحدودي: هل التخطيط سليم على مقاس شاشة الهاتف؟ هل تلتف البطاقات إلى سطر جديد عند ضيق النافذة؟ الطبقة الثالثة، التحقق المعماري: هل الكود مقروء؟ هل CSS مضمّن داخل وسم style واحد؟ هل توجد تكرارات واضحة؟

إذا أردت أيضا تنفيذ وظيفة إحدى بطاقات الأدوات (وليس المدخل فقط)، يمكنك إضافتها تدريجيا — أضف وظيفة أداة واحدة في كل مرة، وتحقق منها قبل إضافة التالية. هذا هو التطبيق العملي لاستراتيجية “الخطوات الصغيرة السريعة”.

\booksection{ثانياً: المشروع الثاني — نظام تسجيل المصروفات المتكامل (مستوى متقدم)}

\booksubsection{التفكير في المتطلبات وكتابة PRD}

الهدف: تطبيق تسجيل مصروفات متكامل (Full-Stack) بواجهة أمامية وخادم خلفي وقاعدة بيانات. هذا المشروع أعقد بكثير من الأول، ويحتاج إلى منهجية هندسة المتطلبات الواردة في الفصل 40.

اكتب وثيقة PRD أولا. وفق قالب القسم الثالث من الفصل 40، اكتب نظرة عامة على المشروع وقائمة الوظائف والحل التقني ونموذج البيانات وخطة التطوير. وفي اختيار الحل التقني احترم مبدأ Simplicity First (البساطة أولا) — واجهة أمامية بـ React دون إدارة حالات معقدة مثل Redux، وخادم خلفي بـ Node.js دون خدمات مصغرة (Microservices)، وقاعدة بيانات SQLite بدلا من PostgreSQL في مرحلة التطوير.

تتفرع خطة التطوير إلى 5 خطوات: الخطوة 1 الهيكل الأساسي لواجهة API الخلفية (خادم Express + SQLite + واجهات CRUD)؛ الخطوة 2 الصفحات الأساسية للواجهة الأمامية (قائمة السجلات + نموذج الإضافة)؛ الخطوة 3 الربط بين الواجهة والخادم (استدعاءات API + عرض البيانات)؛ الخطوة 4 وظائف التصنيف والتصفية؛ الخطوة 5 الملخص الشهري ورسم بياني بسيط. وبعد كل خطوة يكون لديك “شيء قابل للاستخدام”.

\booksubsection{إعداد ملف ذاكرة المشروع}

قبل البدء بالبرمجة، أنشئ ملف ذاكرة المشروع أولا. هذه هي الممارسة الجوهرية لهندسة السياق — بحيث يعرف الذكاء الاصطناعي خلفية المشروع عند بداية كل حوار.

\begin{lstlisting}
# المشروع: تطبيق متكامل لتتبع المصروفات
## نظرة عامة
تطبيق لتتبع المصروفات الشخصية يدعم تسجيل الإيرادات والمصروفات والتصنيف والملخصات الشهرية.
## الحزمة التقنية
- الواجهة الأمامية: React + TypeScript + Vite
- الواجهة الخلفية: Node.js + Express + TypeScript
- قاعدة البيانات: SQLite (عبر better-sqlite3)
- أنواع مشتركة بين الواجهتين الأمامية والخلفية
## اتفاقيات الـ API
- المسار الأساسي: /api
- صيغة الاستجابة: { success: boolean, data: any, error: string }
- أسلوب RESTful
## نموذج البيانات
Record { id, amount, type, category, note, date }
## قواعد الكود
- وضع Strict في TypeScript
- مسافة بادئة بمسافتين
- دوال مع تعليق JSDoc
- التعليقات ورسائل الـ commit بالعربية
\end{lstlisting}

\booksubsection{التنفيذ خطوة بخطوة}

اتبع خطة التطوير، وأنجز كل خطوة في حوار مستقل. وفيما يلي مثال على تعليمات الخطوة الأولى:

\begin{lstlisting}
أحتاج إلى إنشاء هيكل واجهة الـ API الخلفية.
المتطلبات:
1. خادم Express، المنفذ 3001
2. قاعدة بيانات SQLite، اسم الجدول records
3. حقول جدول records:
   id (TEXT PRIMARY KEY)
   amount (REAL)
   type (TEXT: 'income'|'expense')
   category (TEXT)
   note (TEXT)
   date (TEXT: YYYY-MM-DD)
4. نقاط CRUD:
   GET /api/records - جلب كل السجلات
   POST /api/records - إضافة سجل
   PUT /api/records/:id - تحديث سجل
   DELETE /api/records/:id - حذف سجل
5. وسيط معالجة الأخطاء
أعطني البنية الكاملة لكود الواجهة الخلفية.
\end{lstlisting}

لاحظ خصائص هذه التعليمات: تصف الخطوة الحالية وحدها (الهيكل الخلفي) دون خلط متطلبات الواجهة الأمامية؛ وتحدد الحقول والواجهات بوضوح (قيود كافية للذكاء الاصطناعي) وتترك التنفيذ التفصيلي للذكاء الاصطناعي؛ وتلتزم بمبدأ Goal-Driven Execution — قول “ماذا نفعل” لا “كيف ننفذها”.

بعد أن يعطيك الذكاء الاصطناعي الكود، تحقق فورا: شغل الخادم واختبر كل واجهة بـ curl أو Postman للتأكد من عملها. وبعد التأكد افتح حوارا جديدا للخطوة 2 (الصفحات الأساسية للواجهة الأمامية). كل خطوة تتبع دورة “تعليمات — تحقق — حوار جديد”.

\booksection{ثالثاً: المشروع الثالث — تطبيق معزز بالذكاء الاصطناعي (مستوى عال)}

\booksubsection{تصور المشروع}

الهدف: تطبيق يتصل بواجهة النموذج الكبير API — مثل “مساعد تعلم بالذكاء الاصطناعي” يدخل المستخدم فيه محتوى التعلم، فيولد الذكاء الاصطناعي ملخصا ويربط نقاط المعرفة ذات الصلة ويوصي بمصادر التعلم. تميز هذا المشروع بأن تطبيقك نفسه يستدعي الذكاء الاصطناعي، وستحتاج إلى معالجة قضايا جديدة مثل استدعاء API وهندسة الموجهات ومعالجة الأخطاء.

أهم درس يعلّمك إياه هذا المشروع: عندما “تستخدم الذكاء الاصطناعي لكتابة كود يستدعي الذكاء الاصطناعي”، عليك الانتباه بشكل خاص لهندسة الموجهات. فالموجه الذي يستخدمه تطبيقك عند استدعاء واجهة النموذج الكبير API يجب أن يلتزم هو الآخر بمبادئ هذا الدليل — الوضوح والإيجاز والتوجه نحو الهدف.

\booksubsection{النقاط التقنية الأساسية}

يتضمن هذا المشروع عدة نقاط تقنية لم تواجهها في المشروعين السابقين:

\begin{itemize}
  \item إدارة مفاتيح API: لا يجوز كتابة المفتاح داخل الكود مطلقا، وتدار عبر متغيرات البيئة؛
  \item معالجة الاستجابات المتدفقة (Streaming): ردود النماذج الكبيرة متدفقة عادة، وتحتاج إلى عرضها كلمة بكلمة؛
  \item التحكم في استهلاك Token: تحتسب النماذج الكبيرة حسب الـToken، ويلزم التحكم في الاستهلاك على مستوى التطبيق؛
  \item إدارة قوالب الموجهات: لكل وظيفة قالب موجه مختلف، ويلزم إدارتها بشكل منظم.
\end{itemize}

يمكنك إسناد كل هذه النقاط التقنية إلى الذكاء الاصطناعي، لكن عليك فهم “لماذا” وراء كل نقطة — لماذا لا يجوز كتابة المفتاح داخل الكود؟ لماذا نحتاج الاستجابات المتدفقة؟ لماذا يجب التحكم في استهلاك الـToken؟ حين تفهم “اللماذا” تستطيع الحكم على معقولية حل الذكاء الاصطناعي.

\booksection{رابعاً: قائمة سير العمل العملي العامة}

أي كان حجم المشروع، إليك قائمة سير عمل عامة للبرمجة العملية بالذكاء الاصطناعي. راجعها في كل مرة تبدأ فيها مشروعا جديدا للتأكد من عدم إغفال أي خطوة أساسية.

\begin{center}
\begin{tabularx}{\linewidth}{|X|X|X|}
\hline
\textbf{المرحلة} & \textbf{ماذا نفعل} & \textbf{المبدأ المقابل} \tabularnewline
\hline
1. التفكير في المتطلبات & راجع عناصر المتطلبات الخمسة وحدد بوضوح ماذا تريد أن تفعل & Think Before Coding \tabularnewline
2. تفكيك المتطلبات & قسم المتطلبات الكبيرة إلى خطوات صغيرة قابلة للتحقق المستقل & Simplicity First \tabularnewline
3. كتابة PRD & نظم المتطلبات في وثيقة منظمة تستخدمها مرجعا & Goal-Driven \tabularnewline
4. إعداد ملف الذاكرة & أنشئ ملف CLAUDE.md أو ما يشابهه & Harness Engineering \tabularnewline
5. التنفيذ خطوة بخطوة & حوار لكل خطوة: تعليمات — تحقق — الخطوة التالية & Surgical Changes \tabularnewline
6. الاختبار والتحقق & اختبر صحة الوظائف فورا بعد كل خطوة & Iterate \& Verify \tabularnewline
7. مراجعة الكود & فحص معقولية البنية والأمان وقابلية الصيانة & Iterate \& Verify \tabularnewline
8. التكرار والتحسين & تحسين تكراري بناء على نتائج الاختبار والمراجعة & Iterate \& Verify \tabularnewline
\hline
\end{tabularx}
\end{center}

\begin{center}\small\color{gray}قائمة سير العمل العملي العامة للبرمجة بالذكاء الاصطناعي\end{center}

\begin{tcbpractice}{لنطبق عملياً: أتمم "المشروع الأول" هذا الأسبوع}
اجعل “موقع الأدوات الشخصي” الوارد في القسم الأول من الفصل 44 واجبك الأول في البرمجة بالذكاء الاصطناعي: خذ 5 دقائق لمراجعة العناصر الخمسة، وأرسل التعليمات إلى الذكاء الاصطناعي (أي أداة تفي بالغرض)، ثم اقبل العمل وفق التحقق الثلاثي. أتممه خلال 30 دقيقة — ويولد أول أعمالك في البرمجة بالذكاء الاصطناعي. ولا تنس حفظه بعد الإنجاز، فهو تذكار انتقالك من “قارئ” إلى “ممارس”.

\end{tcbpractice}

\begin{tcbsummary}{خلاصة الفصل}
\begin{itemize}
  \item المشروع الأول (تأسيسي): موقع أدوات بواجهة أمامية خالصة — تجربة الدورة الكاملة، ينجز في 30 دقيقة
  \item المشروع الثاني (متقدم): نظام تسجيل مصروفات متكامل — PRD + ملف ذاكرة + تنفيذ في 5 خطوات، وكل خطوة "قابلة للاستخدام"
  \item المشروع الثالث (عال): مساعد تعلم بالذكاء الاصطناعي — استخدام الذكاء الاصطناعي لكتابة كود يستدعي الذكاء الاصطناعي، مع الانتباه لهندسة الموجهات
  \item سير العمل العام في 8 خطوات: التفكير في المتطلبات التفكيك PRD ملف الذاكرة التنفيذ خطوة بخطوة الاختبار المراجعة التكرار
  \item كل خطوة تلتزم دورة "تعليمات — تحقق — حوار جديد"
\end{itemize}

\end{tcbsummary}
